\documentclass[12pt]{article}
\usepackage[T1]{fontenc}
\usepackage[utf8]{inputenc}
\usepackage{lmodern}
\usepackage{textcomp}
\usepackage{enumitem}
\usepackage{geometry}
\geometry{margin=1in}
\begin{document}
\begin{center}
Answer Key - Health Sciences - K-12 Level Assessment 10 \
\small (Answers are ordered exactly as in the question paper.)
\end{center}

\section*{Multiple Choice}
\begin{enumerate}[label=\arabic*.]
\item \textbf{Answer:} D
\\ \textit{Options:} A) Low protein and low GI; B) High protein and high GI; C) Low protein and high GI; D) High protein and low GI
\\ \textit{Note:} A high protein and low glycemic index (GI) diet is recommended for weight maintenance because it promotes satiety, helps regulate blood sugar levels, and reduces cravings, ultimately aiding in the control of calorie intake.
\item \textbf{Answer:} D
\\ \textit{Options:} A) Folate; B) Niacin; C) Riboflavin; D) Thiamin
\\ \textit{Note:} Thiamin, also known as vitamin B$_{1}$, is essential for the formation of the coenzyme thiamine pyrophosphate, which is necessary for the oxidative decarboxylation of pyruvate in the process of converting carbohydrates into energy.
\item \textbf{Answer:} A
\\ \textit{Options:} A) Consumers with phenylketonuria must avoid the consumption of the sweetener aspartame; B) Consumers with phenylketonuria must avoid the consumption of the sweetener saccharin; C) Consumers with phenylketonuria must avoid the consumption of the sweetener sucralose; D) Consumers with phenylketonuria must avoid the consumption of the sweetener acesulfame K
\\ \textit{Note:} Consumers with phenylketonuria must avoid aspartame because it contains phenylalanine, an amino acid that individuals with this condition cannot metabolize, leading to harmful health effects.
\item \textbf{Answer:} C
\\ \textit{Options:} A) glucose; B) Alpha limit dextrins; C) Sucrose; D) polysaccharides
\\ \textit{Note:} Oral bacteria synthesize extracellular glucans such as dextran and mutan from sucrose, as sucrose serves as a key substrate that bacteria can ferment to produce these polysaccharides, which contribute to dental plaque formation.
\item \textbf{Answer:} D
\\ \textit{Options:} A) Vitamin D; B) Vitamin C; C) Vitamin A; D) Folate, vitamins B$_{6}$ and B$_{12}$
\\ \textit{Note:} Folate, vitamin B$_{6}$, and vitamin B$_{12}$ are essential for the metabolism of homocysteine, as they are involved in its conversion into other beneficial compounds, thereby helping to lower its levels in the bloodstream.
\end{enumerate}

\section*{True / False}
\begin{enumerate}[label=\arabic*.]
\setcounter{enumi}{5}
\item \textbf{Answer:} True
\\ \textit{Options:} A) True; B) False
\\ \textit{Note:} As incomes rise and urbanization occurs, people tend to adopt more Western dietary patterns that include higher consumption of meat and dairy, driven by increased availability, affordability, and changing cultural preferences.
\item \textbf{Answer:} True
\\ \textit{Options:} A) True; B) False
\\ \textit{Note:} The statement is true because congenital disorders that impair leptin secretion disrupt the regulation of appetite and energy balance, leading to obesity, abnormal growth, hypothyroidism, and hyperinsulinaemia due to the hormone's critical role in these metabolic processes.
\item \textbf{Answer:} True
\\ \textit{Options:} A) True; B) False
\\ \textit{Note:} Adaptive thermogenesis is true because it refers to the body's ability to alter its heat production in response to variations in calorie intake and environmental temperature, helping to maintain energy balance and regulate body temperature.
\item \textbf{Answer:} False
\\ \textit{Options:} A) True; B) False
\\ \textit{Note:} Thiamin deficiency in alcoholics primarily leads to Wernicke-Korsakoff syndrome, not Cushing syndrome, which is associated with cortisol levels.
\item \textbf{Answer:} True
\\ \textit{Options:} A) True; B) False
\\ \textit{Note:} Plasma lipoproteins contain a hydrophobic core made up of triacylglycerols and cholesterol esters, which allows them to transport lipids through the aqueous environment of the bloodstream.
\end{enumerate}

\section*{Long Form Response}
\begin{enumerate}[label=\arabic*.]
\setcounter{enumi}{10}
\item \textbf{Answer:} Intestinal "brakes" and "accelerators" play crucial roles in regulating digestion by controlling the speed at which food moves through the digestive system. The "brakes" slow down motility to allow for nutrient absorption, while the "accelerators" increase motility to promote faster digestion and movement of food. Dietary starch is often misunderstood; it is actually a strong stimulator of small intestinal motility because it triggers the release of insulin from the pancreas. This insulin release aids in the digestion and absorption of glucose, contradicting the notion that starch would slow down the digestive process. Thus, the statement regarding dietary starch and insulin might be considered false because it overlooks the stimulatory effect that starch has on intestinal activity.
\\ \textit{Note:} correct as it accurately describes the opposing functions of intestinal "brakes" and "accelerators" in digestion and clarifies the misconception about dietary starch being a stimulant for insulin release, emphasizing its role in enhancing small intestinal motility instead.
\item \textbf{Answer:} Mitochondria and peroxisomes both play important roles in fatty acid ß-oxidation, which is the process by which fatty acids are broken down for energy. While peroxisomes do participate in this process, they are not the major site for fatty acid ß-oxidation; that role belongs to the mitochondria. Mitochondria contain the necessary enzymes and are specialized for efficiently converting fatty acids into energy through ß-oxidation. Peroxisomes primarily assist in the initial steps of breaking down very long-chain fatty acids before sending them to mitochondria for further processing. Therefore, it is incorrect to say that peroxisomes are the major site of fatty acid ß-oxidation, as mitochondria are the main location where this energy-producing process occurs.
\\ \textit{Note:} correct because it accurately identifies mitochondria as the primary site for fatty acid ß-oxidation due to their specialized enzymes and efficiency in energy conversion, while acknowledging that peroxisomes are involved only in the initial steps of the process.
\end{enumerate}

\vfill
\noindent\textit{End of Answer Key}
\end{document}